\textbf{Giải thích ví dụ 1}

Ta có thể xác định rằng sau 3 lần cân đầu tiên, đồng xu $1$ là đồng xu nặng nhất, 
nhưng không thể biết điều đó chỉ sau 2 lần cân. 
Do đó, số nguyên đầu tiên trong kết quả đúng là $3$.

Tương tự, ta có thể xác định rằng đồng xu $2$ là đồng xu nhẹ nhất sau 4 lần cân, 
nhưng không phải sau 3 lần cân. 
Vì vậy, số nguyên thứ hai trong kết quả đúng là $4$.

Lưu ý rằng cả hai thứ tự $2,4,3,1$ và $2,3,4,1$ đều là các cách sắp xếp hợp lệ của trọng lượng các đồng xu.  
Do đó, đồng xu $3$ có thể có cấp bậc $2$ hoặc $3$, đều phù hợp với tất cả các kết quả cân, 
vì vậy cấp bậc của đồng xu $3$ không bao giờ được xác định chắc chắn.  
Tương tự, đồng xu $4$ cũng không bao giờ có cấp bậc xác định.

Nếu kết quả của bạn là \texttt{1 1 -1 -1}, bạn sẽ đạt được 50\% số điểm cho test nhỏ này.
